% !TEX program = xelatex
% 컴파일: xelatex repro_report.tex  (또는 lualatex).  한글은 kotex 패키지 사용.
\documentclass[11pt]{article}

\usepackage{kotex}          % 한국어 (xelatex/lualatex 권장; pdflatex도 가능)
% 시스템에 설치된 Noto CJK KR 폰트를 명시 지정 (xelatex/lualatex + tectonic)
\setmainhangulfont{Noto Serif CJK KR}
\setsanshangulfont{Noto Sans CJK KR}
\setmonohangulfont{Noto Sans CJK KR}
\usepackage[margin=25mm]{geometry}
\usepackage{booktabs}
\usepackage{array}
\usepackage{longtable}
\usepackage{multirow}
\usepackage{xcolor}
\usepackage{enumitem}
\usepackage{amsmath}
\usepackage[hidelinks]{hyperref}
\usepackage{titlesec}

\newcommand{\yes}{\textcolor{green!55!black}{\textbf{O}}}
\newcommand{\no}{\textcolor{red!70!black}{\textbf{X}}}
\newcommand{\warn}{\textcolor{orange!85!black}{\textbf{$\triangle$}}}
\newcommand{\code}[1]{\texttt{\small #1}}

\titleformat{\section}{\large\bfseries}{\thesection}{0.6em}{}
\setlist{itemsep=2pt, topsep=3pt}

\title{\textbf{official\_matrix\_88 원본 충실 재현 보고서}\\[2pt]
\large 커널 생성 에이전트 $\times$ 공식 벤치마크 매트릭스: 현황과 완전 재현 계획}
\author{}
\date{2026년 7월 13일}

\begin{document}
\maketitle

\begin{abstract}
\noindent
본 보고서는 \code{official\_matrix\_88}(11개 커널 생성 에이전트 $\times$ 8개 공식 벤치마크 = 88셀)을
각 에이전트의 \textbf{실제 구현/모델}로 재현하고 각 벤치마크의 \textbf{공식 오라클}로 채점하는
``원본 충실(faithful) 재현''의 현황과 남은 완전 재현 계획을 정리한다.
기존에 저장소에 커밋되어 있던 88셀 결과는 11개 에이전트가 모두 하나의 제너릭 드라이버를 통해
``smoke test용 단순 torch 코드''를 생성한 것으로, 에이전트의 실제 성능을 나타내지 않는다.
본 계획의 대상 범위는 \textbf{실제 구현/모델로 충실 재현이 가능한 8개 에이전트}로 한정한다:
cudaforge, autokernel, autotriton, kernelllm, ksearch, kernelskill, drkernel, incoder32b.
모델이 비공개인 2종(cuda\_l1, cuda\_agent)과 AMD 전용 1종(geak)은 범위에서 제외한다.
따라서 대상 매트릭스는 \textbf{8개 에이전트 $\times$ 8개 벤치마크 $=$ 64셀이며 전 셀이 원본 충실(faithful)}이다.
\textbf{현재까지의 진척:} (1) \textbf{8개 faithful 에이전트를 전부 실제 구현/모델로 재현}하여
KernelBench(level1/19\_ReLU)의 공식 오라클 판정을 확보하였고(6/8 correct, 2/8은 모델 자체의
충실한 실패), (2) \textbf{8개 벤치마크의 공식 오라클을 전부 이 Blackwell(sm\_120) 박스에서 배선}하여
각각 구조화된 공식 판정을 반환함을 확인하였다. 즉 완전 매트릭스를 위한 두 축(에이전트 축, 벤치마크 축)이
모두 준비되었다. 이하에서 재현 결과, 오라클 배선 현황, GPU 구성, 시간/비용, 결론, 논문 연결을 제시한다.
\end{abstract}

\section{배경}
\code{official\_matrix\_88}은 서베이 논문의 실증(empirical) 섹션을 위한 자기완결형 하네스로,
11개 오픈소스 커널 생성 에이전트를 8개 벤치마크의 \emph{공식 오라클}로 통일 평가한다.
그러나 저장소에 커밋된 결과는 모든 에이전트가 \code{drivers/generic\_llm\_kernel\_driver.py}
하나를 경유했고, 그 프롬프트가 명시적으로
``\emph{for this smoke test, do not use Triton/CUDA/inline extensions}''라고 지시하여
실제 커널이 아닌 단순 torch 코드를 생성하였다. 따라서 커밋된 88셀은
\textbf{파이프라인 동작 확인(smoke)}일 뿐 에이전트 성능이 아니다.

본 작업의 목표는 각 에이전트를 \textbf{원저자의 실제 구현/모델 그대로} 구동하여
``실제 에이전트 생성 $\rightarrow$ 공식 오라클 판정''의 진짜 데이터를 얻는 것이다.
모든 에이전트는 로컬 백본(예: Qwen2.5-Coder-14B) 또는 각자의 공개 모델을 사용하며,
평가는 벤치마크의 공식 오라클(예: KernelBench는 \code{telemetry/instrumented\_final\_eval.py})로 수행한다.

\section{재현 방식의 3가지 분류}
에이전트를 실제로 조사한 결과, ``에이전트''는 세 가지 서로 다른 형태를 가진다.
\begin{itemize}
  \item \textbf{라이브 에이전트(프레임워크 내장 루프)}: 생성$\to$컴파일/벤치$\to$피드백$\to$개선을
        스스로 수행. 예: cudaforge(\code{main.py}, NCU 하드웨어 피드백 10라운드),
        kernelskill(KernelMem, 메모리 루프), ksearch, drkernel(KernelGYM).
  \item \textbf{모델-온리 릴리스}: 저장소에 에이전트 코드가 없고 \emph{파인튜닝된 모델}이 산출물.
        모델에 프롬프트하면 커널을 생성. 예: autotriton(8B), kernelllm(8B), incoder32b(32B).
  \item \textbf{아티팩트 replay}: RL 학습 모델이 비공개이고, 사전 생성된 최적화 코드만 공개.
        라이브 실행 불가, 공개 코드를 replay하여 채점만 가능. 예: cuda\_l1, cuda\_agent.
\end{itemize}

\section{재현 결과: faithful 8종 전부 (KernelBench level1/19\_ReLU)}
각 에이전트의 \textbf{실제 구현/모델}로 재현하고 공식 KernelBench 오라클
(\code{telemetry/instrumented\_final\_eval.py})로 채점한 결과다.
평가 하드웨어는 단일 NVIDIA RTX PRO 6000 (Blackwell, sm\_120)이다.

\begin{table}[h]
\centering
\begin{tabular}{@{}lllll@{}}
\toprule
\textbf{에이전트} & \textbf{유형} & \textbf{재현 방식(원본 그대로)} & \textbf{공식 판정} & \textbf{speedup} \\
\midrule
cudaforge   & 라이브 CUDA   & \code{main.py} 3라운드 + NCU 하드웨어 피드백 & \yes\ correct & 0.70$\times$ \\
autokernel  & 라이브 CUDA   & 네이티브 브리지 + Qwen keep/revert 루프    & \yes\ correct & 1.008$\times$ \\
autotriton  & 모델 Triton(8B) & AutoTriton RL 모델 서빙                 & \yes\ correct & 0.99$\times$ \\
kernelllm   & 모델 Triton(8B) & KernelLLM 공식 프롬프트 템플릿           & \warn\ compiled, \no\ incorrect & --- \\
ksearch     & 라이브 Triton & world-model 탐색 2라운드                 & \yes\ correct & 0.994$\times$ \\
kernelskill & 라이브 CUDA   & KernelMem 메모리 루프 + 자가수정          & \yes\ correct & 0.695$\times$ \\
drkernel    & 모델 Triton(14B) & Dr.Kernel-14B reasoning 생성           & \yes\ correct & 0.994$\times$ \\
incoder32b  & 모델 CUDA(32B) & IndustrialCoder(vLLM Llama 오버라이드)   & \no\ compiled=False & --- \\
\bottomrule
\end{tabular}
\caption{faithful 8종의 공식 오라클 판정(단일 태스크 ReLU). \textbf{6/8 correct, 2/8 충실한 실패}
(kernelllm은 차원을 잘못 하드코딩, incoder32b는 CUDA 디바이스 코드를 C++ 컴파일 유닛에 배치 --- 둘 다
파이프라인 문제가 아니라 모델 자체의 오류). ReLU는 메모리 바운드라 speedup이 1.0 부근에 수렴하므로,
이는 성능 순위가 아니라 파이프라인/정확도 검증이다.}
\end{table}

\noindent 이는 기존 smoke 결과와 질적으로 다른, 논문에 사용 가능한 실측 데이터다.
예컨대 cudaforge는 실제로 NCU를 실행해 하드웨어 메트릭을 피드백하여 커널을
0.5030$\to$0.5311(내부 점수)로 개선하였고, ksearch는 world-model 탐색으로 1라운드 실패 후
2라운드에서 통과 커널을 찾았으며, autotriton/drkernel/kernelllm은 각자의 RL 모델로 커널을 생성하였다.
2건의 실패도 ``correct$<$eval은 정상적 벤치 결과''(README)로서 정직한 데이터포인트다.

\section{에이전트별 완전 재현 가능성}
11개 에이전트를 실제 구현/모델로 재현 가능한지 재분류하였다.

\begin{table}[h]
\centering
\begin{tabular}{@{}lll@{}}
\toprule
\textbf{에이전트} & \textbf{정체} & \textbf{재현 가능성} \\
\midrule
cudaforge   & 라이브(CUDA, NCU)        & \yes\ faithful \\
autokernel  & 라이브(CUDA/Triton 브리지) & \yes\ faithful \\
autotriton  & 모델-온리(8B, Triton)     & \yes\ faithful \\
kernelllm   & 모델-온리(8B, Triton)     & \yes\ faithful \\
ksearch     & 라이브(frontier LLM)      & \yes\ faithful \\
kernelskill & 라이브(KernelMem 메모리)   & \yes\ faithful \\
drkernel    & RL 모델(14B)+gym          & \yes\ faithful \\
incoder32b  & 모델-온리(\textbf{32B})    & \warn\ faithful (무거움) \\
cuda\_l1    & RL 모델 비공개            & \warn\ replay만 \\
cuda\_agent & RL 모델 비공개(데이터셋만)  & \warn\ replay/대체만 \\
geak        & \textbf{AMD MI 전용}      & \no\ NVIDIA 재현 불가 \\
\bottomrule
\end{tabular}
\caption{재현 가능성 재분류. geak(AMD 전용)는 제외, cuda\_l1/cuda\_agent는 모델 비공개로 replay만 가능.}
\end{table}

\noindent\textbf{핵심:} 11개 전부를 원본 충실 재현하는 것은 불가능하다. 3개 에이전트가 근본적으로 막힌다
(모델 비공개 2: cuda\_l1, cuda\_agent / AMD 전용 1: geak). 이는 본 작업의 한계가 아니라
원저자의 미공개 또는 타 하드웨어 요구에 기인한다. 따라서 본 계획은 이 3종을 범위에서 제외하고,
\textbf{충실 재현이 가능한 8종만}을 대상으로 하여 매트릭스 전 셀을 원본 충실로 유지한다.

\section{재현 범위: faithful 8종 (64셀 전부 충실)}
모델 비공개 2종과 AMD 전용 1종을 제외하고, 실제 구현/모델로 충실 재현이 가능한
\textbf{8 에이전트 $\times$ 8 벤치마크 = 64셀}만을 대상으로 한다. 전 셀이 원본 충실이므로
replay/N-A 라벨이 필요 없다.
\begin{itemize}
  \item \textbf{faithful 8종}: cudaforge, autokernel, autotriton, kernelllm, ksearch, kernelskill, drkernel, incoder32b
  \item \textbf{제외 3종}: cuda\_l1(모델 비공개), cuda\_agent(모델 비공개), geak(AMD 전용)
  \item 8 벤치마크: kernelbench, robust\_kbench, tritonbench\_t, tritonbench\_g, multikernelbench, backendbench, pareval, sol\_execbench
\end{itemize}
$\Rightarrow$ \textbf{64셀, 전부 faithful.}

\section{벤치마크 축: 8개 공식 오라클 배선 완료 (Blackwell)}
매트릭스의 다른 축인 8개 벤치마크의 \textbf{공식 오라클(채점기)}을 전부 이 Blackwell 박스에서 실행하여,
각각 구조화된 공식 판정을 반환함을 확인하였다. 러너(\code{official\_all\_matrix\_v1.py})의
\code{evaluate\_*} 디스패치는 \textbf{코드 수정 없이} 동작했으며, 필요한 것은 환경 설정
(\code{TORCH\_CUDA\_ARCH\_LIST=12.0}, 실행별 새 \code{TORCH\_EXTENSIONS\_DIR}, 오프라인 플래그,
LLM 엔드포인트)과 \code{prepare}(태스크 매니페스트 생성)뿐이었다.

\begin{table}[h]
\centering
\begin{tabular}{@{}llc@{}}
\toprule
\textbf{벤치마크} & \textbf{공식 오라클(채점기)} & \textbf{배선 검증(cudaforge, $T{=}1$)} \\
\midrule
kernelbench      & \code{instrumented\_final\_eval.py}         & \yes\ (8종 채점에 사용) \\
robust\_kbench   & kb instrumented (\code{bench.sh matrix})   & \yes\ official\_eval=1, correct=1 \\
tritonbench\_t   & TritonBench EVAL (call/exe/eff)            & \yes\ official\_eval=1, correct=1 \\
tritonbench\_g   & TritonBench EVAL                           & \yes\ official\_eval=1, correct=1 \\
multikernelbench & \code{eval\_single\_runner.py}              & \yes\ official\_eval=1, correct=1 \\
backendbench     & BackendBench CLI (directory backend)      & \yes\ official\_eval=4/8, ``not covered'' 정상 \\
pareval          & ParEval C++ 드라이버 (\code{run-all.py})    & \yes\ official\_eval=1, correct=0 \\
sol\_execbench   & \code{python -m sol\_execbench.cli.main}    & \yes\ official\_eval=1, correct=1 \\
\bottomrule
\end{tabular}
\caption{8개 벤치마크 오라클 전부 Blackwell에서 정상 작동(구조화된 공식 판정 반환).
smoke 후보(cudaforge 제너릭 드라이버)로 배선을 검증한 것이므로 correct 값 자체는 참고용이며,
핵심은 각 오라클이 이 하드웨어에서 실행되어 판정을 낸다는 점이다. backendbench의 add/mul/sub/div는
로컬 테스트 스위트가 없어 ``not covered''로 정상 처리된다(README 문서화 동작).
prepare로 생성된 태스크 수: kernelbench 249, robust 200, tritonbench\_t 166, \_g 184,
multikernel 401, backendbench 15, pareval 60, sol 21.}
\end{table}

\noindent\textbf{요약:} 에이전트 축(8종 재현)과 벤치마크 축(8 오라클 배선)이 모두 준비되어,
완전 매트릭스(64셀) 실행에 필요한 인프라가 완성되었다.

\section{GPU 구성 (4장, 최상위)}
\textbf{추천: 4$\times$ NVIDIA B200 (180GB, Blackwell)} --- 32B 모델도 여유, sm\_120 검증 완료.
저렴한 대안: 4$\times$ H100 80GB 또는 H200 141GB.
\begin{itemize}
  \item GPU 0--1: 에이전트 모델 서빙(vLLM). 에이전트마다 백본이 달라 모델 스왑 필요
        (Qwen14B, AutoTriton8B, KernelLLM8B, drkernel14B, incoder32B $\dots$).
  \item GPU 2--3: 커널 컴파일 및 공식 오라클 eval 병렬 실행.
  \item 병렬 효율 약 70\%(모델 스왑, 라운드 내 직렬성) $\Rightarrow$ 유효 배속 $\approx 2.8\times$.
\end{itemize}

\section{시간 및 비용 추정}
셀당 시간 $=$ 태스크 수 $T \times$ (생성 $+$ eval). 로컬 실측 기반 태스크당 평균:

\begin{table}[h]
\centering
\begin{tabular}{@{}lccc@{}}
\toprule
\textbf{에이전트 유형} & \textbf{생성/태스크} & \textbf{eval/태스크} & \textbf{합} \\
\midrule
반복형 ($\times 5$: cudaforge, autokernel, ksearch, kernelskill, drkernel) & $\sim$6분 & $\sim$2분 & 8분 \\
단발형 ($\times 3$: autotriton, kernelllm, incoder32b)                      & $\sim$2분 & $\sim$2분 & 4분 \\
\bottomrule
\end{tabular}
\caption{태스크당 평균 소요(로컬 RTX PRO 6000 실측 기반). 벤치마크 1개(8 에이전트)당
$= 5\times8 + 3\times4 = 52T$분, $\times 8$벤치 $= 416T$분(단일 GPU 직렬).}
\end{table}

4-GPU 유효배속($\div 2.8$) 적용 시:

\begin{table}[h]
\centering
\begin{tabular}{@{}ccccc@{}}
\toprule
\textbf{$T$ (셀당 태스크)} & \textbf{벽시계(4$\times$GPU)} & \textbf{GPU-시간} & \textbf{4$\times$B200(\textasciitilde\$20/hr)} & \textbf{4$\times$H100(\textasciitilde\$10/hr)} \\
\midrule
1 (현재 수준)        & $\sim$2.5시간 & $\sim$10   & \textasciitilde\$50    & \textasciitilde\$25 \\
10 (권장 최소)       & $\sim$25시간  & $\sim$99   & \textasciitilde\$500   & \textasciitilde\$250 \\
50                   & $\sim$5.2일   & $\sim$496  & \textasciitilde\$2{,}480 & \textasciitilde\$1{,}240 \\
250 (KernelBench 전량) & $\sim$26일   & $\sim$2{,}480 & \textasciitilde\$12{,}400 & \textasciitilde\$6{,}200 \\
\bottomrule
\end{tabular}
\caption{균일 $T$ 가정의 상한 근사. 실제 $T$는 벤치마크별로 다르다(pareval/sol 등은 태스크 수가 적음).}
\end{table}

\noindent\textbf{일회성 통합 엔지니어링은 이미 완료되었다.} CudaForge/autokernel 등을 붙이는 데
CUDA OOM, torch 확장 stale lock, \code{PYTHONPATH}, trial timeout, 커스텀 아키텍처 로딩 등 다수의
디버깅이 필요했으나(부록 참조), \textbf{8종 에이전트 통합과 8개 오라클 배선이 모두 끝난 상태}이다.
따라서 남은 것은 위 표의 \textbf{GPU 실행 비용($T$에 비례)뿐}이다.
\begin{itemize}
  \item 에이전트 통합(8종): \textbf{완료} --- 드라이버 및 서빙 절차 확립.
  \item 벤치마크 오라클 배선(8종): \textbf{완료} --- 전부 Blackwell에서 검증.
  \item 남은 작업 = smoke 제너릭 드라이버 대신 \textbf{실제 에이전트를 8개 벤치마크 전부에 $T{>}1$로 실행}.
\end{itemize}
\textbf{권장 운영점: $T=10$}. 4$\times$H100 기준 \textbf{\textasciitilde 25시간 / \textasciitilde \$250}이면
통계적 최소 의미를 확보하며, 통합/배선이 완료되었으므로 \textbf{바로 실행 가능}하다.

\section{결과물}
\textbf{8$\times$8 = 64셀 매트릭스(전부 faithful)} --- 셀마다 공식 오라클 지표를 담는다.
\begin{itemize}
  \item 정확도: \code{correct\_rate}, \code{pass@1}
  \item 성능: \code{speedup} / \code{fast\_p} (fast\_1.0/1.5/2.0/3.0)
  \item 원출력 및 verdict JSON 전량 보존(재현성).
\end{itemize}
즉 ``에이전트 $\times$ 벤치마크''를 동일 공식 오라클로 사과-대-사과 비교한 표와 벤치별 난이도 프로파일.
모든 셀이 원본 충실이므로 replay/N-A 각주가 필요 없다.

\section{내릴 수 있는 결론}
\begin{enumerate}
  \item \textbf{에이전트 간 공정 비교}: 서로 다른 프레임워크를 동일 오라클/하드웨어에서 재평가하여
        벤치마크 계열별로 어느 에이전트가 실제로 정확하고 빠른 커널을 내는지 규명.
  \item \textbf{벤치마크 난이도 구조}: kb-계열 vs operator-fill(backendbench) vs 병렬(pareval) 등의 성공률 패턴.
  \item \textbf{패러다임 효과}: 반복형(NCU 피드백)/메모리형/단발 모델/RL 모델의 상대 우위.
  \item \textbf{정직한 경계}: 대상은 충실 재현이 가능한 8종이며, 모델 비공개 2종/AMD 1종은
        범위 밖임을 명시하여 결론 범위를 한정.
  \item 기존 smoke 대비 ``진짜 커널 생성'' 재평가라는 점.
\end{enumerate}
\noindent\emph{못 내리는 결론}: geak/AMD 비교, cuda\_agent/cuda\_l1의 라이브 성능(모델 비공개),
$T$가 작을 때의 세밀한 순위.

\section{논문 연결}
\begin{itemize}
  \item \textbf{Table 2 (taxonomy)}: 에이전트 $\times$ 벤치마크 분류에 재현 방식(faithful/replay/N-A) 표기 추가.
  \item \textbf{Contribution 2 (인프라)}: 모든 벤치마크의 공식 오라클을 하나로 배선한 이식형 통합 하네스
        및 에이전트별 실제 구현 통합(\code{repro\_env.sh}, 각 드라이버). \textbf{8종 에이전트 재현과
        8개 오라클 배선을 Blackwell(sm\_120)에서 end-to-end로 검증 완료.}
  \item \textbf{Main Results (Table 3)}: 8$\times$8 = 64셀 완전 충실 매트릭스(정확도 + fast\_p).
  \item \textbf{Analysis}: 위 결론 1--3.
  \item \textbf{Limitations / Threats to Validity}: 범위 제외 3종(모델 비공개 2: cuda\_l1, cuda\_agent /
        AMD 전용 1: geak), controlled vs native backbone, $T$ 규모, 하드웨어(Blackwell sm\_120) 특수성.
  \item \textbf{Reproducibility 부록}: 아래 하드-원 이슈 및 커맨드.
\end{itemize}

\section{부록: 재현 환경 하드-원 이슈}
Blackwell RTX PRO 6000 상에서 KernelBench 형식 eval + vLLM 서빙 시 확인된 비자명 이슈와 해법:
\begin{enumerate}[label=A\arabic*.]
  \item \textbf{아키텍처 핀}: \code{TORCH\_CUDA\_ARCH\_LIST=12.0}(sm\_120) 미설정 시 6개 아키텍처를
        모두 컴파일하여 3분+ 소요 $\to$ eval 타임아웃. 핀 설정 시 콜드 컴파일 $\sim$33초.
  \item \textbf{torch 확장 stale lock}: 타임아웃으로 SIGKILL된 컴파일이
        \code{\textasciitilde/.cache/torch\_extensions/.../lock}(무확장자)를 남겨 이후 프로세스가 무한 대기.
        실행마다 새 \code{TORCH\_EXTENSIONS\_DIR} 사용으로 회피.
  \item \textbf{GPU 공유}: vLLM과 eval이 한 GPU를 공유. \code{--gpu-memory-utilization 0.4}로 서빙하여
        eval에 여유 확보(0.85는 OOM 유발).
  \item \textbf{vLLM flashinfer 불일치}: \code{FLASHINFER\_DISABLE\_VERSION\_CHECK=1}로 우회.
  \item \textbf{autokernel 통합}: 플레이북이 \code{from kernels.cuda.\_compile import compile\_cuda}를
        사용하므로 \code{PYTHONPATH}에 autokernel 루트 추가 필요. 콜드 \code{compile\_cuda}가 30초
        trial 타임아웃 내부에서 실행되므로 \code{AK\_TRIAL\_TIMEOUT} 상향 필요.
  \item \textbf{NCU}: 본 장비에서 sudo 없이 GPU 카운터 접근 가능 $\to$ cudaforge 하드웨어 피드백 정상.
  \item \textbf{kernelllm}: 출력이 torch-inductor 스타일이며 torch$\ge$2.11에서 제거된
        \code{grid} 심볼을 import한다 $\to$ 커널 로직은 그대로 두고 \code{grid}만 복원하는
        버전-호환 심을 앞에 덧댐.
  \item \textbf{kernelskill(KernelMem)}: 릴리스에 \code{bench\_ref\_inputs\_0.py} 템플릿 누락
        $\to$ 부모 프로젝트 CudaForge의 동일 파일 이식. NCU 경로가 \code{/root/...}로 하드코딩
        (무시하고 진행). 프롬프트가 컴파일 플래그를 A100/sm\_80으로 하드코딩.
  \item \textbf{ksearch(K-Search)}: KernelBench를 \code{kernelbench} 패키지로 import
        $\to$ \code{third\_party/KernelBench/src}를 \code{PYTHONPATH}에 추가.
  \item \textbf{incoder32b(IndustrialCoder)}: 커스텀 아키텍처 \code{IQuestCoderForCausalLM}은
        실제로는 Llama 계열(QKV bias 없음, RMSNorm, SwiGLU, GQA, rope\_theta 500000)이다.
        transformers \code{trust\_remote\_code} 경로는 rope 초기화 비호환으로 깨진 출력을 냈고,
        vLLM \code{--hf-overrides architectures=LlamaForCausalLM}로 서빙하니 정상 로드/생성되었다.
  \item \textbf{벤치마크 오라클(8종)}: 러너의 \code{evaluate\_*} 배선은 코드 수정 없이 Blackwell에서
        동작하였고, 위 A1(아키텍처 핀)/A2(stale lock) 환경 설정과 \code{prepare}만으로 전부 실행되었다.
\end{enumerate}

\end{document}
